% J3.4: condiciones de contorno de la viga y sus primeras raíces de la ecuación característica.
\documentclass[tikz,border=4pt]{standalone}
\usepackage{resonancia}
\begin{document}
\begin{tikzpicture}
  \tikzset{viga/.style={line width=2.2pt, color=rAzul}}
  % 1. empotrada-libre
  \begin{scope}[shift={(0,0)}]
    \fill[apoyo] (-0.3,-0.35) rectangle (0,0.35); \draw[trazo] (0,-0.35) -- (0,0.35);
    \draw[viga] (0,0) -- (2.6,0);
    \node[etiqueta, anchor=north, align=center] at (1.3,-0.5) {\textbf{empotrada–libre}\\[-1pt]$\cos x+1/\cosh x=0$\\[-1pt]$x_1=1{,}875$};
  \end{scope}
  % 2. biapoyada
  \begin{scope}[shift={(4.2,0)}]
    \draw[viga] (0,0) -- (2.6,0);
    \draw[trazo] (0,0) -- (-0.2,-0.3) -- (0.2,-0.3) -- cycle;
    \draw[trazo] (2.6,0) -- (2.4,-0.3) -- (2.8,-0.3) -- cycle; \draw[trazo] (2.3,-0.38) -- (2.9,-0.38);
    \node[etiqueta, anchor=north, align=center] at (1.3,-0.5) {\textbf{biapoyada}\\[-1pt]$\sin x=0$\\[-1pt]$x_1=\pi=3{,}142$};
  \end{scope}
  % 3. empotrada-apoyada
  \begin{scope}[shift={(8.4,0)}]
    \fill[apoyo] (-0.3,-0.35) rectangle (0,0.35); \draw[trazo] (0,-0.35) -- (0,0.35);
    \draw[viga] (0,0) -- (2.6,0);
    \draw[trazo] (2.6,0) -- (2.4,-0.3) -- (2.8,-0.3) -- cycle; \draw[trazo] (2.3,-0.38) -- (2.9,-0.38);
    \node[etiqueta, anchor=north, align=center] at (1.3,-0.5) {\textbf{empotrada–apoyada}\\[-1pt]$\sin x-\cos x\tanh x=0$\\[-1pt]$x_1=3{,}927$};
  \end{scope}
  % 4. libre-libre
  \begin{scope}[shift={(12.6,0)}]
    \draw[viga] (0,0) -- (2.6,0);
    \node[etiqueta, anchor=north, align=center] at (1.3,-0.5) {\textbf{libre–libre}\\[-1pt]$\cos x-1/\cosh x=0$\\[-1pt]$x_1=4{,}730$};
  \end{scope}
  \node[etiqueta, color=rGris, anchor=north] at (8.0,-1.95) {$f_n=\dfrac{x_n^{2}}{2\pi L^{2}}\sqrt{\dfrac{EI}{\rho A}}$, con $x_n=\beta_n L$ la raíz $n$ de la ecuación característica};
\end{tikzpicture}
\end{document}
